\documentclass[12pt]{article}
\usepackage[utf8]{inputenc}
\usepackage[T1]{fontenc}
\usepackage{hyperref}
\title{Warduji Language Resources for K-12 Curriculum Development}
\author{Compiled by Claude Sonnet 5.5 (claude-sonnet-5-5) for Andr\'as Kornai}
\date{}
\begin{document}
\maketitle

\section{Introduction}
Warduji (code wrd; Glottocode ward1243) does not appear to be a documented language. The ISO 639-3 Registration Authority retired the code wrd (Change Request 2021-015, effective 2022-01-20), and the change form named no replacement. Glottolog still lists the entry under a ``Bookkeeping'' grouping, meaning a catalogue placeholder and not a language family. Global Recordings Network still shows wrd with a population figure of 6,780, but this was not independently verified. Searches found no grammar, dictionary, primer, school programme or community organisation for it.

\section{Primary Educational Resources}
None found.

\section{Linguistic and Reference Materials}
\begin{itemize}
\item Glottolog record of the ISO 639-3 Registration Authority 2021 change: \url{https://glottolog.org/resource/reference/id/641369}
\item Glottolog page for Warduji (ward1243): \url{https://glottolog.org/resource/languoid/id/ward1243.bigmap.html}
\item ISO 639-3 change request form for wrd: \url{https://mtb.jdhuang.people.aws.dev/assignments/lang-subtags-templates/wrd.txt}
\end{itemize}

\section{Religious and Community Resources}
Global Recordings Network has a page for wrd; its content was not examined: \url{https://globalrecordings.net/ja/language/wrd}

\section{Digital and Audio Resources}
See the Global Recordings Network page above. No other audio, video or software was found.

\section{Curriculum Development Context}
A curriculum cannot be planned for this code, because it is not clear that a distinct language exists. The first step would be to establish what, if anything, the code referred to.

\end{document}